\specialchapter{BIBLIOGRAPHY}

\begin{enumerate}
\def\labelenumi{(\Roman{enumi})}
\item
  J. HESSEL, Krysallometrie oder Krystallonomie und Krystallographie (1830, repr. in Ostwald's Klassiker, nos. 88 and 89, Leipzig (Teubner), 1897)
\item
  A. Bravais, Mémoires sur les polyèdres de forme symétrique, Journal de Math., (1), v. XIV (1849), p.~141-180.
\item
  A. MÖBIUS: a) Ueber das Gesetz der Symmetrie der Krystalle und die Anwendung dieses Gesetze auf die Eintheilung der Krystalle in Systeme, J. für die reine und angewandte Math., v. XLIII (1852), p.~365-374 (= Gesammelte Werke, v. II, Leipzig (Hirzel), 1886, p.~349-360); b) Theorie der symmetrischen Figuren, Gesammelte Werke, v. II, Leipzig (Hirzel), 1886, p.~561-708.
\item
  L. SCHLÄFLI, Theorie der vielfachen Kontinuität. Denkschr. der Schweiz. naturforsch. Gesellschaft, v. XXXVIII (1901), p.~1-237 (= Ges. math. Abhandlungen, v. I, Basel (Birkhäuser), 1950, p.~167-387).
\item
  W. R. HAMILTON, Memorandum respecting a new system of roots of unity, Phil. Mag., (4), v. XII (1856), p.~446.
\item
  C. JORDAN, Mémoire sur les groupes de mouvements, Ann. di Mat., v. Il (1868-69), p.~167-215 and 322-345 (= Oeuvres, v. IV, Paris (Gauthier-Villars), 1964, p.~231-302).
\item
  E. GOURSAT, Sur les substitutions orthogonales et les divisions régulières de l'espace, Ann. Ec. Norm. Sup., (3), v. VI (1889), p.~9-102.
\item
  W. KILLING, Die Zusammensetzung der stetigen endlichen Transformationsgruppen: I) Math. Ann., v. XXXI (1888), p.~252-290; II) ibid., v. XXXIII (1889), p.~1-48; III) ibid., v. XXXIV (1889), p.~57-122; IV) ibid., v. XXXVI (1890), p.~161-189.
\item
  E. CARTAN: a) Sur la structure des groupes de transformations finis et continus (Thesis). Paris (Nony), 1894 (= Oeuvres complètes, Paris (Gauthier-Villars), 1952, v. I₁, p.~137-287); b) Sur la réduction à sa forme canonique de la structure d'un groupe de transformations fini et continu, Amer. Journ. of Math., v. XVIII (1896), p.~1-46 (= Oeuvres complètes, Paris (Gauthier-Villars), 1952, v. I₁, p.~293-353); c) Le principe de dualité et la théorie des groupes simples et semi-simples, Bull. Sci. Math., v. XLIX (1925), p.~361-374 (= Oeuvres complètes, Paris (Gauthier-Villars), 1952, v. I₁, p.~555-568); d) La géométrie des groupes simples, Ann. di Mat., (4), v. IV (1927), p.~209-256 (= Oeuvres complètes, Paris (Gauthier-Villars), 1952, v. I₂, p.~793-840); e) Sur certaines formes riemanniennes remarquables des géométries à groupe fondamental simple, Ann. Ec. Norm. Sup., (3), v. XLIV (1927), p.~345-467 (= Oeuvres complètes, Paris (Gauthier-Villars), 1952, v. I₂, p.~867-989); f) Complément au mémoire ``Sur la géométrie des groupes simples'', Ann. di Mat., v. V (1928), p.
\end{enumerate}

253-260 (= Oeuvres complètes, Paris (Gauthier-Villars), 1952, v. I₂, p.~1003-1010).

\begin{enumerate}
\def\labelenumi{(\Alph{enumi})}
\setcounter{enumi}{23}
\tightlist
\item
  R. FRICKE and F. KLEIN, Theorie der automorphen Funktionen, Leipzig (Teubner), 1897.
\end{enumerate}

\begin{enumerate}
\def\labelenumi{(\Roman{enumi})}
\setcounter{enumi}{10}
\item
  L. E. DICKSON: a) Theory of linear groups in an arbitrary field, Trans. Amer. Math. Soc., v. II (1901), p.~363-394; b) A new system of simple groups, Math. Ann., t. LX (1905), p.~137-150; c) A class of groups in an arbitrary realm connected with the configuration of the 27 lines on a cubic surface, Quart. Journ. of Math., v. XXXIII (1901), p.~145-173 and v. XXXIX (1908), p.~205-209.
\item
  A. SPEISER. Theorie der Gruppen von endlicher Ordnung, Berlin (Springer), 1924.
\item
  H. WEYL, Theorie der Darstellung kontinuerlicher halb-einfacher Gruppen durch lineare Transformationen, Math. Zeitschr., v. XXIII (1925), p.~271-309, v. XXIV (1926), p.~328-395 and 789-791 (= Selecta, Basel-Stuttgart (Birkhäuser), 1956, p.~262-366).
\item
  H. S. M. COXETER: a) Groups whose fundamental regions are simplexes, Journ. Lond. Math. Soc., v. VI (1931), p.~132-136; b) The polytopes with regular prismatic figures, Proc. Lond. Math. Soc., (2), v. XXXIV (1932), p.~126-189; c) Discrete groups generated by reflections, Ann. of Math., (2), v. XXXV (1934), p.~588-621; d) The complete enumeration of finite groups of the form \(\mathbf{R}_i^2 = (\mathbf{R}_i\cdot \mathbf{R}_j)^{k_{ij}} = 1\) , Journ. Lond. Math. Soc., v. X (1935), p.~21-25; e) Regular polytopes, New York (Macmillan), 1948 (2nd ed., 1963); f) The product of generators of a finite group generated by reflections, Duke Math. Journ., v. XVIII (1951), p.~765-782.
\end{enumerate}

(XIV bis) H. S. M. COXETER and H. WEYL, The structure and representations of continuous groups (Inst. for Adv. Study, mimeographed notes by N. Jacobson and R. Brauer, 1934-35): Appendix.

\begin{enumerate}
\def\labelenumi{(\Roman{enumi})}
\setcounter{enumi}{14}
\item
  H. S. M. COXETER and W. O. J. MOSER, Generators and relations for discrete groups, Ergebn. der Math., Neue Folge, Bd. 14, Berlin (Springer), 1957 (2nd ed., 1963).
\item
  B. L. van der WAERDEN, Die Klassification der einfachen Lieschen Gruppen, Math. Zeitschr., v. XXXVII (1933). p.~446-462.
\item
  E. WITT, Spiegelungsgruppen und Aufzählung halbeinfacher Lieschen Ringe, Abh. Math. Sem. Hamb. Univ., v. XIV (1941), p.~289-322.
\item
  E. STIFEL, Ueber eine Beziehung zwischen geschlossenen Lie'sche Gruppen und diskontinuierlichen Bewegungsgruppen euklidischer Räume und ihre Anwendung auf die Aufzählung der einfachen Lie'schen Gruppen, Comm. Math. Helv., v. XIV (1941-42), p.~350-380.
\item
  C. CHEVALLEY: a) Sur la classification des algèbres de Lie simples et de leurs représentations, C. R. Acad. Sci., v. CCXXVII (1948), p.~1136-1138; b) Invariants of finite groups generated by reflections, Amer. Journ. of Math., v. LXXVII (1955), p.~778-782; c) Sur certains groupes simples, Tohôku Math. Journ., (2), v. VII (1955), p.~14-66; d) Classification des groupes de Lie algébriques, 2 vols., Paris (Inst. H. Poincaré), 1956-58.
\item
  HARISH CHANDRA: a) On some applications of the universal enveloping algebra of a semi-simple Lie algebra, Trans. Amer. Math. Soc., v. LXX (1951), p.~28-96; b) On a lemma of Bruhat, Journ. de Math., (9), v. XXXV (1956), p.~203-210.
\item
  F. BRUHAT, Représentations induites des groupes de Lie semi-simples complexes, C. R. Acad. Sci., v. CCXXXVIII (1954), p.~437-439.
\item
  A. BOREL, Groupes linéaires algébriques, Ann. of Math., (2), v. LXIV (1956). p.~20-80.
\item
  A. J. COLEMAN, The Betti numbers of the simple groups, Can. Journ. of Math.. v. X (1958). p.~349-356.
\item
  R. STEINBERG. Finite reflection groups. Trans. Amer. Math. Soc., v. XCI (1959). p.~493-504.
\item
  J. TITS: a) Groupes simples et géométries associées, Proc. Int. Congress Math., Stockholm, 1962, p.~197-221; b) Théorème de Bruhat et sous-groupes paraboliques. C. R. Acad. Sci., v. CCLIV (1962), p.~2910-2912; c) Algebraic and abstract simple groups. Ann. of Math., (2), v. LXXX (1964), p.~313-329.
\item
  N. IWAHORI and H. MATSUMOTO, On some Bruhat decomposition and the structure of the Hecke rings of \(p\) -adic Chevalley groups, Publ. Math. I.H.E.S., no. 25 (1965), p.~5-48.
\item
  A. BOREL and J. TITS. Groupes réductifs. Publ. Math. I.H.E.S., no. 27 (1965), p.~55-150.
\item
  F. BRUHAT and J. TITS. Groupes algébriques simples sur un corps local, Proc. Conf. on local Fields, p.~23-36. Berlin (Springer), 1967.
\item
  R. CARTER, Simple groups and simple Lie algebras, Journ. Lond. Math. Soc., v. XL (1965), p.~193-240.
\end{enumerate}

